\documentclass[11pt,a4paper]{article}

% ------------------------------------------------------------
% Pacotes
% ------------------------------------------------------------
\usepackage[utf8]{inputenc}
\usepackage[T1]{fontenc}
\usepackage[brazil]{babel}
\usepackage{lmodern}
\usepackage{microtype}
\usepackage{amsmath,amssymb}
\usepackage{enumitem}
\usepackage{geometry}
\usepackage{titlesec}
\usepackage{hyperref}

% ------------------------------------------------------------
% Layout
% ------------------------------------------------------------
\geometry{
    top=1.8cm,
    bottom=1.8cm,
    left=1.9cm,
    right=1.9cm
}

\setlength{\parindent}{0pt}
\setlength{\parskip}{3pt}

\setlist[itemize]{
    leftmargin=1.35em,
    itemsep=0.8pt,
    topsep=2pt,
    parsep=0pt
}

\setlist[enumerate]{
    leftmargin=1.55em,
    itemsep=0.8pt,
    topsep=2pt,
    parsep=0pt
}

\titleformat{\section}
    {\bfseries\normalsize}
    {}
    {0pt}
    {}

\titlespacing*{\section}
    {0pt}
    {7pt}
    {2pt}

\hypersetup{
    colorlinks=true,
    linkcolor=black,
    urlcolor=blue
}

% ------------------------------------------------------------
% Documento
% ------------------------------------------------------------
\begin{document}

\begin{center}
    {\large\bfseries MO438 -- Redes Complexas}\\[2pt]
    \textbf{Proposta de Projeto}
\end{center}

\vspace{-2mm}

% ============================================================
\section*{1. Título}

Da rede lexical à rede contextual: reorganização das representações
internas de Transformers com Redes Complexas

% ============================================================
\section*{2. Breve descrição do projeto}

Transformers representam tokens por vetores em espaços de alta dimensão. Na
entrada do modelo, ocorrências de um mesmo token possuem o mesmo
\textit{embedding} lexical. À medida que atravessam as camadas do Transformer,
essas representações incorporam informações do contexto e podem se reorganizar
no espaço vetorial.

Este projeto propõe representar essa transformação como uma sequência de
redes complexas. Os vértices corresponderão a tokens ou ocorrências de
tokens, e as arestas conectarão representações semelhantes. A comparação das
redes permitirá investigar em que medida a estrutura permanece organizada pela
identidade lexical e em que medida passa a refletir o
contexto ao longo das camadas.

A hipótese central não é que a identidade lexical desapareça, mas que exista
uma reorganização progressiva das vizinhanças e comunidades, cuja intensidade
possa ser quantificada por ferramentas de Redes Complexas.

% ============================================================
\section*{3. Objetivo geral do projeto}

Caracterizar, por meio de redes complexas, como a estrutura de
similaridade entre representações de tokens se transforma entre o espaço
lexical de entrada e as representações contextuais produzidas pelas diferentes
camadas de um Transformer.

% ============================================================
\section*{4. Objetivos específicos do projeto}

\begin{itemize}

    \item Construir uma \textbf{rede lexical de tipos},
    conectando tokens com embeddings de entrada
    semelhantes.

    \item Construir uma \textbf{rede lexical por ocorrência} e redes
    contextuais com aproximadamente $1.5\times10^4$ ocorrências, mantendo os
    mesmos vértices em todas as redes para permitir comparação direta.

    \item Quantificar quanto a \textbf{identidade lexical} determina a
    vizinhança de cada ocorrência em diferentes camadas.

    \item Verificar a hipótese de que as vizinhanças se reorganizam ao longo
    das camadas à medida que as representações incorporam informação
    contextual.

    \item Identificar e comparar comunidades nas diferentes redes,
    verificando se sua organização permanece associada aos tipos lexicais ou
    passa a refletir diferentes contextos.

    \item Analisar tokens com múltiplas ocorrências, especialmente aqueles
    encontrados em contextos distintos, buscando identificar separações
    estruturais entre suas representações.

    \item Calcular e comparar as principais propriedades das redes:
    $|V|$, $|E|$, grau médio, distribuição de graus, densidade,
    clusterização global e local, distância média, diâmetro, componentes
    conexos e distribuição de seus tamanhos.

    \item Visualizar a evolução da estrutura das redes ao longo das camadas do
    Transformer.

\end{itemize}

% ============================================================
\section*{5. Origem dos dados}

Será utilizada uma amostra de textos reais em português, inicialmente obtida da 
Wikipédia. Os textos serão processados utilizando o tokenizer do
próprio Transformer analisado.

A análise principal utilizará aproximadamente $1.5\times10^4$ ocorrências de
tokens. A frequência por tipo será controlada para impedir que tokens muito
frequentes dominem artificialmente a topologia. Tokens especiais serão
excluídos e, quando apropriado, a análise será concentrada em tokens de
conteúdo.

% ============================================================
\section*{6. Como os objetivos serão atingidos?}

Será utilizado um Transformer aberto, preferencialmente entre
4B e 8B parâmetros. Para cada ocorrência $i$, serão extraídos seu
embedding lexical $e_{t_i}$ e os estados internos $h_i^{(\ell)}$ de camadas
representativas do início, meio e fim do modelo.

Serão construídos três tipos de rede:

\begin{enumerate}

    \item \textbf{Rede lexical de tipos:} um vértice para cada tipo de token,
    representado por seu embedding de entrada $e_t$.

    \item \textbf{Rede lexical por ocorrência:} os vértices são as ocorrências
    analisadas, mas cada ocorrência $i$ é representada por
    $x_i=e_{t_i}$. Assim, ocorrências de um mesmo token possuem inicialmente
    representações idênticas.

    \item \textbf{Redes contextuais:} os mesmos vértices são representados
    pelos estados $x_i=h_i^{(\ell)}$ em cada camada $\ell$ analisada.

\end{enumerate}

A similaridade entre dois vértices será medida pelo cosseno,

\[
s_{ij} =
\frac{x_i^\top x_j}
{\lVert x_i\rVert\,\lVert x_j\rVert},
\]

e cada vértice será conectado aos seus $k$ vizinhos mais próximos, formando
redes k-NN esparsas. O mesmo $k$ e a mesma regra de simetrização serão
utilizados em todas as camadas. Diferentes valores de $k$ serão testados como
análise de robustez.

A reorganização das vizinhanças será medida pela similaridade de Jaccard entre
os conjuntos de vizinhos de uma mesma ocorrência em diferentes camadas. A
dominância lexical será medida pela fração de vizinhos de uma
ocorrência que correspondem ao mesmo tipo de token.

% ============================================================
\section*{7. Como os resultados serão avaliados?}

O cumprimento dos objetivos será avaliado em três níveis. Primeiro, será
verificado se todas as redes propostas puderam ser construídas com escala
superior a $10^4$ vértices e se as métricas estruturais requeridas pela
disciplina foram calculadas.

Segundo, serão comparadas quantitativamente as redes entre camadas. A evolução
da dominância lexical indicará quanto a identidade do token continua
determinando sua vizinhança, enquanto a similaridade de Jaccard permitirá medir
a intensidade da reorganização local. As comunidades e métricas globais
mostrarão se essa transformação também ocorre em escalas maiores da rede.

Por fim, serão analisados tokens individuais com múltiplas ocorrências para
verificar se representações inicialmente equivalentes passam a ocupar regiões
distintas da rede em função do contexto. Os resultados serão testados para
diferentes valores de $k$ e limites de frequência.

% ============================================================
\section*{8. Planejamento e cronograma}

\begin{itemize}

    \item \textbf{Semanas 1--2:} seleção do modelo e corpus, definição da
    amostragem e extração dos embeddings e estados internos.

    \item \textbf{Semanas 2--3:} construção da rede lexical de tipos, da rede
    lexical por ocorrência e das redes contextuais.

    \item \textbf{Semanas 3--4:} cálculo das métricas estruturais, persistência
    de vizinhança e dominância lexical.

    \item \textbf{Semanas 4--5:} detecção de comunidades, análise intra-token,
    visualizações e testes de robustez.

    \item \textbf{Semanas 5--6:} consolidação dos resultados e redação do
    relatório final.

\end{itemize}

\end{document}
